% Transcription — fonds Alexandre Grothendieck, shelfmark 13, batch 3 (pages 41-51).
% Established from the facsimile archives/batches/13/batch-03.pdf (a mirror of
% https://grothendieck.umontpellier.fr/13.pdf), per the skill
% transcribe-grothendieck. The batch file carries no cover sheet and is cut from
% the folder starting at page 41: PDF page 1 is archivists' page 41, PDF page N
% is archivists' page N+40. Checked against the pencilled numbers bottom-left —
% "45", "46", "47", "48" and "50" read clearly on the corresponding leaves and
% fix the alignment across the whole batch. This is the folder's last batch.
%
% Pass: Opus 5 (claude-opus-5), 2026-09-19 — first pass, unchecked against the pages by a human.
%
% What the batch is. Eleven leaves in his hand, of which ten carry mathematics
% (page 44 is blank and gets no \page{}). They fall into three runs:
%
%   41-43   the close of the run batch 2 was in the middle of: the description
%           of the motivic Galois group G of a base S together with the torsor
%           of periods P, the equivariant coverings, and the Hodge datum
%           f : \widetilde{P}_{C} -> \Gamma_{G}; it ends on page 43 with a
%           five-item summary A) to E) of the whole structure and a N.B. asking
%           what part of it determines the rest.
%   45-49   a new, separately headed piece, « Foncteurs pleinement fidèles sur
%           catégories des motifs, liés à théorie de Hodge », numbered 1) to 6):
%           six candidate realisation functors on motives, each asked whether it
%           is fully faithful, and under which conjecture.
%   50-51   a short closing note, « Remarques sur la caractérisation des classes
%           de coh. D.R. essentiellement algébriques », numbered 1°) to 3°).
%
% Where the run stood at page 41. Batch 2's last leaves (38-40) set up the
% situation afresh: S regular of finite type over Q, k the algebraic closure of
% Q in Gamma(S, O_S) with an embedding k -> C, M the category of motives over S,
% G = Aut^{tensor}(T) the motivic Galois group, P = Isom^{tensor}(T_{B}, T_{DR})
% a torsor under G_{S} carrying an integrable connection, and then the analytic
% objects \overline{S}_{C} (universal cover of S(C) at xi_0), \widetilde{P}_{C}
% and the map f into the space \Gamma_{G} of Hodge structures on G_{R}, with its
% conditions a), b), c). Page 41 opens in the middle of that list of conditions
% and simply continues it: nothing is restated.
%
% The argument, page by page:
%
%   41   having fixed xi_0 in S(C), the equivariant coverings are equivalent to
%        a homomorphism w : pi_1(S(C), xi_0) -> G(Q) inside G(C) together with a
%        point x_0 of P_{xi_0}; this gives the unique q : \overline{S}_{C} ->
%        P_{C}^{an} with q(\bar{s}\sigma) = q(\bar{s})w(\sigma), hence
%        \widetilde{P}_{C} = \overline{S}_{C} x^{\Gamma} G(Q) inside P_{C}^{an},
%        a reduction of the structure group from G(C) to G(Q). Then the map
%        f : \widetilde{P}_{C} -> \Gamma_{G} and its conditions a) holomorphy,
%        b) equivariance f(xg) = g^{-1}f(x), c) compatibility at each point s.
%   42   the same data read on \overline{S}_{C} rather than \widetilde{P}_{C};
%        then the adelic layer: the covering \widetilde{P}_{C} x^{G(Q)} G(A),
%        A the finite adeles, and the factorisation of Gamma -> G(Q) -> G(A)
%        through the algebraic fundamental group pi_1(S, xi_0); the quotient
%        \widetilde{P}_{C}/G^{0}(Q) identified with (P/G^{0})_{C}^{an}.
%   43   the summary A) to E) — reductive group G over Q; principal bundle P
%        with integrable connection; reduction of the structure group to G(Q);
%        the holomorphic equivariant f; the adelic restriction — followed by a
%        N.B.: A) B) C) plus the single value f(xi_0) should in principle
%        reconstruct D) entirely and the data coming from E). Closes on two
%        questions, an analytic continuation formula for f and the relation
%        between D) and E), the second glossed « Jugendtraum ? ».
%   45   § 1): S a regular connected Q-scheme with a geometric point
%        xi : Spec(C) -> S; xi^{*} : Mot(S) -> Mot(C) is faithful, and one
%        studies the composite (T_{B})_{xi} into Mod_{f}(Q). To a motive M one
%        attaches the triple {T_{DR}(M_y), T_{B}(M_xi), phi_M} with T_{DR}
%        filtered (a), or the triple {T_{Hdg}(M), T_{B}(M_xi), psi_M} with
%        T_{Hdg} bigraded (b); both functors are fully faithful granted the
%        Hodge conjecture. The circled marginals name the linear-algebra data
%        each triple amounts to: « réseau et vectoriel filtré », « réseau et
%        vectoriel bigradué ».
%   46   § 2) « Via Tate »: the pair {T_{B}(M_xi), rho_{l,M}} with rho_{l,M}
%        the l-adic representation of pi = pi_1(S, xi); fully faithful when S is
%        of finite type over Q, granting Tate and using a result of Griffiths.
%        Over C one takes instead pi' = pi_1(S(C), xi) and rho_M, provided T_B
%        is replaced by T_{B-H}, the Hodge structure. § 3) the quadruple
%        (T_{DR}(M), sigma, T_{B}(M_y), phi) with sigma the Gauss-Manin
%        stratification.
%   47   § 4) « Double niveau »: S = Spec k, k algebraic over Q, the system
%        (4.1) {T_{DR}(M), T_{B}(M), phi} with neither filtration nor Hodge
%        structure. Reduced to k = \overline{Q}, the question becomes whether
%        (4.2) {T_{B}(M), sigma} is fully faithful, sigma and T_{DR}(M) inside
%        T_{B}(M) tensor C determining each other. Brought inside the Galois
%        group: a class \dot{g} in G(C)/G(\overline{Q}), and for a G-module V
%        the question reduces to whether V^{G} = V \cap g(\overline{V}).
%   48   the same question worked out: x in \overline{V} \cap g(\overline{V})
%        expressed through the stabiliser G_x, a subgroup \overline{T} of
%        \overline{G}, the image \overline{R} of \overline{T} x \overline{T}
%        under (s,t) -> s^{-1}t, and a subgroup H of G, ending on a Borel
%        subgroup. The page is written over itself and heavily cancelled.
%        § 5) begins: over a k of finite type over Q the functor is no longer
%        fully faithful, and he breaks off with « idiot ! » before hoping for a
%        disproportion result between T_B and T_{DR}.
%   49   without hypothesis on k beyond being algebraically closed, (4.2) should
%        be fully faithful; this would follow from the double-level conjecture
%        of 4) together with Tate or Hodge, plus § 6): for S of finite type the
%        functor of 3) without the filtration is fully faithful.
%   50   « Remarques sur la caractérisation des classes de coh. D.R.
%        essentiellement algébriques » : K a subfield of C, M a motive over K;
%        T_{DR}(M) is contained a) in T_{DR}(M)^{00}_{C} and b) in
%        (T_{DR}(M)_{C})_{R}. Is there a converse? 1°) not for a) and b)
%        together; 2°) the converse to a) alone fails once the motivic Galois
%        group has dimension at least 3, through V' = P x^{G_K} V and a
%        representation for which V'^{T} differs from V'^{G'}.
%   51   3°) what remains open: a converse to b), taking all the embeddings
%        K -> C at once and, if need be, conjoining condition a). Three lines,
%        and the folder stops.
%
% Notation fixed for the batch, and held to. S is the base, xi (or xi_0) its
% geometric point, y its generic point — he writes M_{y} and M_{xi} side by side
% on page 45 and the distinction is his. \overline{S}_{C} is the universal cover
% of S(C), \widetilde{P} and \widetilde{P}_{C} the principal covering with
% structure group G(Q), P_{C}^{an} = P(C) the analytic torsor; \Gamma is
% pi_1(S(C), xi_0) and \Gamma_{G} the space of Hodge structures on G_{R} — the
% two Gammas are his and are not the same letter twice by accident, the second
% always carrying the subscript G. G is the motivic Galois group, G^{0} its
% neutral component, G_x the stabiliser of x; A is the ring of finite adeles,
% which he defines in a circled marginal on page 42. T_{B}, T_{DR}, T_{Hdg} are
% the Betti, de Rham and Hodge realisations, T_{B-H} the Betti realisation
% carrying its Hodge structure; Mot(S), Mot(C), Hodge(Q,C), Mod_{f}(Q) are his
% category names, set in \mathrm. phi_M, psi_M are the comparison isomorphisms,
% rho_{l,M} and rho_M the l-adic and topological monodromy representations.
% \overline{V} is V tensor_{Q} \overline{Q}, and T_{DR}(M)^{00} his superscript
% for the part of weight 0 on page 50. Words he underlines are set \emph.
%
% What does not carry. Pages 43, 45, 47 and 51 are written comparatively fair
% and read almost end to end. Pages 41, 42, 46 and 49 are quicker and lose parts
% of their connective prose. Pages 48 and 50 are a fast palimpsest: page 48
% carries a boxed passage written over, struck through and written again, and
% page 50 a marginal cancelled past any recovery — in both the displayed formulas
% carry and the sentences between them very largely do not. The apparatus is
% dense there in consequence, and that density is the record of what the leaves
% support.
%
% The folder's end. The argument does not close: page 51 is three lines opening
% a third question, 3°), which it states and does not answer, and the leaf — and
% with it the folder — simply stops. No letter falls in these eleven pages; the
% 1965 letter the inventory announces is elsewhere in the folder.

%
% Parallel pass, and what was reconciled afterwards. The folder's three batches
% were read concurrently, in three separate contexts, so batches 2 and 3 could
% not read the preceding batch's finished .tex to inherit its notation; each
% read the tail of the preceding facsimile instead. Two seams were reconciled
% in the orchestrating session once all three had landed, and nothing else was
% touched:
%   — batch 1 had set Q, R, C, Z, S, U, G_m in \mathbf and batches 2 and 3 in
%     \mathbb. Batch 1 was converted to \mathbb (205 occurrences), which is
%     both the corpus convention and the better record of the doubled strokes
%     he draws. No reading changed.
%   — the cover on page 20 had become a \section in batch 1 (closing it) and
%     again in batch 2 (opening it). It is now a heading in batch 2 alone,
%     where the run it names actually begins, and a \note in batch 1.
% Neither was a disagreement about what the leaves say.

\documentclass[11pt,a4paper]{article}
% Le préambule commun à toutes les transcriptions du fonds Grothendieck.
%
% Il tient en une idée : **l'incertitude se compose, elle ne se résout pas.**
% Une transcription qui lisse un mot douteux a détruit la seule chose pour
% laquelle on la faisait — un lecteur doit pouvoir savoir, à chaque endroit,
% ce qui était sur le feuillet et ce qui a été ajouté. Les six macros qui
% suivent sont donc l'appareil critique, et non de la décoration.
%
% Les mêmes commandes sont reconnues par `scripts/render.mjs`, qui produit la
% vue de lecture du site. Une macro ajoutée ici doit l'être aussi là-bas,
% faute de quoi le rendu échouera bruyamment — ce qui est voulu : un rendu
% silencieusement faux est pire qu'un rendu absent.

\ProvidesPackage{grothendieck}[2026/08/08 Transcriptions du fonds Grothendieck]

\usepackage{amsmath,amssymb,amsthm}
\usepackage{mathrsfs}
\usepackage{stmaryrd}
% Les diagrammes commutatifs sont partout dans ce fonds : c'est la langue
% même de la géométrie algébrique de Grothendieck, et une transcription qui
% les rendrait en prose ne transcrirait rien.
\usepackage{tikz-cd}
% Français par défaut — c'est la langue des feuillets — mais l'anglais est
% chargé aussi : la traduction et le résumé sont des documents anglais, et la
% typographie française (espaces avant les deux-points) y serait une faute.
% Ces documents basculent par \selectlanguage{english} en tête de corps.
\usepackage[english,french]{babel}
\usepackage{fontspec}
\usepackage{geometry}
\geometry{a4paper,margin=25mm,marginparwidth=28mm}
\usepackage{marginnote}
\usepackage{xcolor}
% ulem, not soul. `soul`'s \st and \ul are built on a character-by-character
% scan that XeLaTeX's font machinery breaks: the document fails with an
% undefined control sequence rather than degrading. ulem does the same two jobs
% and is Unicode-engine safe. `normalem` stops it hijacking \emph, which the
% transcriptions use for Grothendieck's own emphasis.
\usepackage[normalem]{ulem}
\usepackage{hyperref}
\hypersetup{colorlinks=true,linkcolor=black,urlcolor=black,pdfborder={0 0 0}}

% --- Métadonnées du lot -----------------------------------------------------
% Reprises telles quelles par le script de rendu, qui les affiche en tête de
% la vue de lecture. Elles fixent ce que le fichier prétend couvrir : sans
% elles, une transcription flotte sans point d'attache dans un fonds de seize
% mille feuillets.

\newcommand{\folder}[1]{\gdef\grothFolder{#1}}
\newcommand{\batch}[1]{\gdef\grothBatch{#1}}
\newcommand{\leaves}[2]{\gdef\grothFirst{#1}\gdef\grothLast{#2}}
\newcommand{\foldertitle}[1]{\gdef\grothFoldertitle{#1}}
\gdef\grothFolder{?}\gdef\grothBatch{?}\gdef\grothFirst{?}\gdef\grothLast{?}\gdef\grothFoldertitle{}

% --- Le résumé --------------------------------------------------------------
%
% Il ouvre la lecture modernisée, sous le titre « Résumé », et il est composé
% en petit. Ce n'est pas un ornement : c'est le seul passage du document qui
% ne soit pas des mathématiques, et un lecteur doit pouvoir le reconnaître
% avant de l'avoir lu — puis le sauter s'il connaît déjà le sujet, ou s'y
% arrêter s'il ne le connaît pas. Le filet qui le suit marque la reprise du
% corps.

\newenvironment{resume}
  {\small\setlength{\parskip}{0.45em}}
  {\par\medskip\hrule\medskip}

% --- Mots-clés ---------------------------------------------------------------
%
% En anglais, en clôture du résumé de la lecture modernisée : le vocabulaire
% moderne sous lequel chercher ce que les feuillets construisent. Le manifeste
% du site (`npm run manifest`) les extrait pour étiqueter la cote entière —
% cette ligne est la source unique des tags, il n'y a pas de fichier à côté.
\newcommand{\keywords}[1]{\par\smallskip\noindent
  {\footnotesize\textsc{Keywords} --- \textit{#1}}\par}

% --- L'appareil critique ----------------------------------------------------

% Le numéro de feuillet, en marge. C'est l'ancre qui permet de régler un
% désaccord sur une formule en regardant *un* feuillet plutôt qu'une chemise
% de six cents. Le site s'en sert aussi pour tourner les pages du fac-similé
% quand on fait défiler la transcription.
\newcommand{\leaf}[1]{%
  \marginnote{\footnotesize\color{gray!70}\textsf{#1}}%
  \label{leaf:#1}\ignorespaces}

% Illisible : on ne devine pas. Un mot inventé qui se lit comme les autres est
% le pire résultat possible.
\newcommand{\ill}{\textcolor{red!60!black}{[\,\ldots\,]}}

% Lecture douteuse : le mot est donné, et signalé comme incertain.
\newcommand{\uncertain}[1]{\uline{#1}}

% Ajout de l'éditeur — une lettre manquante, un mot sous-entendu.
\newcommand{\add}[1]{\textcolor{blue!50!black}{[#1]}}

% Biffé par Grothendieck lui-même. Ce qu'il a rayé fait partie du document :
% une rature dit souvent où la pensée a changé de direction.
\newcommand{\struck}[1]{\sout{#1}}

% Note du transcripteur, sur l'état matériel du feuillet.
\newcommand{\note}[1]{\par\smallskip{\small\color{gray!60!black}\textit{[#1]}}\par\smallskip}

% Note marginale de Grothendieck, distincte du corps du texte.
\newcommand{\marginal}[1]{\par\smallskip\noindent\hspace{1em}%
  {\small\color{gray!60!black}#1}\par\smallskip}

% --- Titre ------------------------------------------------------------------

\AtBeginDocument{%
  \begin{center}
    {\large\bfseries Fonds Alexandre Grothendieck --- cote n\textsuperscript{o}~\grothFolder}\\[2pt]
    {\itshape\grothFoldertitle}\\[4pt]
    {\small feuillets \grothFirst--\grothLast\ (lot \grothBatch)}\\[2pt]
    {\footnotesize\color{gray!60!black}Transcription établie d'après le fac-similé numérisé
     par l'Université de Montpellier.\\
     Les lectures incertaines sont soulignées, les illisibles notées [\,\ldots\,],
     les ajouts entre crochets.}
  \end{center}
  \vspace{1em}}

\endinput


\folder{13}
\batch{3}
\pages{41}{51}
\dating{1965}
\watermark{Édition de démonstration}
\foldertitle{Groupe de Galois motivique : notes manuscrites (s.d.), lettre (1965).}

\begin{document}

\section{Le torseur des périodes, le revêtement équivariant et la donnée de Hodge}

\note{Le feuillet 41 reprend en cours de route l'énoncé commencé au feuillet 40 :
les objets $S$, $G$, $P$, $\overline{S}_{\mathbb{C}}$, $\widetilde{P}$ et
$\Gamma_{G}$ y sont posés et ne sont pas redits ici. Aucun titre ne figure en
tête du feuillet ; celui-ci est de l'éditeur.}

\page{41}

Lorsqu'on a fixé un point $\xi_{0} \in S(\mathbb{C})$, \add{et qu'on prend
$T = T_{\xi_{0}}$}, alors la donnée des revêtements équivariants
$\overline{S}_{\mathbb{C}}$ et de $\widetilde{P}$ équivalent à la donnée de

\[ w : \pi_{1}(S(\mathbb{C}), \xi_{0}) \longrightarrow G(\mathbb{Q}) \subset G(\mathbb{C}) \]

\note{Entre crochets, d'une écriture plus fine, il justifie ce morphisme :
« \ill{} intégrable de $\widetilde{P}$, qui \ill{} $\overline{S}_{\mathbb{C}}^{\mathrm{an}}$,
d'où $\Gamma = \pi_{1}(S(\mathbb{C}), \xi_{0}) \to G(\mathbb{C})$, \ill{} qu'il
se factorise par $G(\mathbb{Q})$ ».}

et à un \struck{\ill{}} point $x_{0}$ de $P_{\xi_{0}}$, i.e.\ une trivialisation
de $P_{\xi_{0}}$, qui définit alors un unique morphisme $q$ du rev.\ universel
$\overline{S}_{\mathbb{C}}$ de $S_{\mathbb{C}}$ dans
$P_{\mathbb{C}}^{\mathrm{an}} = P(\mathbb{C})$, qui envoie le pt marqué
$\overline{\xi}_{0}$ en $x_{0}$ ; il satisfait
\[ q(\overline{s}\,\sigma) = q(\overline{s}) \cdot w(\sigma) , \quad \overline{s} \in \overline{S}_{\mathbb{C}} , \ \sigma \in \pi_{1}(S(\mathbb{C}), \xi_{0}) , \]
et définit
\[ \widetilde{P}_{\mathbb{C}} = \overline{S}_{\mathbb{C}} \times^{\Gamma} G(\mathbb{Q})
\hookrightarrow P_{\mathbb{C}}^{\mathrm{an}} \]
\add{restriction du groupe structural \ill{} de $G(\mathbb{C})$ à $G(\mathbb{Q})$}.

Comme structure importante, il y a une application
\[ f : \widetilde{P}_{\mathbb{C}} \longrightarrow \Gamma_{G} \]
où $\Gamma_{G} = $ l'ens.\ des structures de Hodge sur $G_{\mathbb{R}}$
(\struck{compatibles} \struck{associées}, i.e.\ \ill{} liées avec un morphisme
$i : \mathbb{G}_{m,\mathbb{R}} \to G_{\mathbb{R}}$ donné). Cette application $f$
satisfait les conditions

\begin{enumerate}
  \item[a)] $f$ holomorphe ;
  \item[b)] $f(x \cdot g) = g^{-1} f(x)$ \quad $g \in G(\mathbb{Q})$, $x \in \widetilde{P}_{\mathbb{C}}$ ;
  \item[c)] $\forall\, x \in \widetilde{P}_{\mathbb{C}}$ sur le point $s$ de $S$,
  la \struck{filtration} structure de \ill{} $T_{\mathbb{C}}$ sur
  \uncertain{$\operatorname{Rep}(\mathbb{C})$}, définie par
  $f(x) : S \to G_{\mathbb{R}}$, est compatible
\end{enumerate}

\page{42}

avec le torseur $P_{s}$ sous $G_{k(s)}$ et la trivialisation
$p(x) \in P_{s}(\mathbb{C})$.

En termes de $w$ et $q : \overline{S}_{\mathbb{C}} \to P_{\mathbb{C}}^{\mathrm{an}}$,
la donnée de $f$ revient à \add{la} donnée de
\[ f : \overline{S}_{\mathbb{C}} \longrightarrow \Gamma_{G} \]
\struck{satisfaisant} \add{a)} \struck{à} les conditions

\begin{enumerate}
  \item[a)] $f$ holomorphe ;
  \item[b)] $f(x \cdot g) = w(g)^{-1} f(x)$, pour \ill{} ;
  \item[c)] \struck{$f(\ill{})$} $\forall\, \overline{s} \in \overline{S}_{\mathbb{C}}$,
  $f(\overline{s}) : S \to G_{\mathbb{R}}$ est compatible avec le
  torseur $P_{s}$ sous $G_{k(s)}$ et la trivialisation $q(x) \in P_{s}(\mathbb{C})$.
\end{enumerate}

Enfin, \struck{pour tout $k$,} \ill{} on se donne une « variation »
\struck{en vue.} $\widetilde{P}_{\mathbb{C}} \times^{G(\mathbb{Q})} G(\mathbb{A})$
sur $\overline{S}_{\mathbb{C}}^{\mathrm{an}}$, ou un revêtement \ill{} de groupe
$G(\mathbb{Q}_{\ell})$ sur $S$.

\marginal{$\mathbb{A} = $ adèles « finis » sur $\mathbb{Q}$}

En termes de \struck{$\xi_{0}$ et} $\Gamma = \pi_{1}(S(\mathbb{C}), \xi_{0})$, on
se donne une factorisation de $\Gamma \to G(\mathbb{Q}) \longrightarrow G(\mathbb{A})$ par
\[ \Gamma \longrightarrow \pi_{1}(S, \xi_{0}) \longrightarrow G(\mathbb{A}) , \]
$\pi_{1}$ désignant ici le groupe fondamental algébrique. Quant au revêtement
$\widetilde{P}_{\mathbb{C}} \times^{G(\mathbb{Q})} (G/G^{\circ})(\mathbb{A})$, il
est \ill{} \struck{$\widetilde{P}_{\mathbb{C}} \times$} associé :
\[ \widetilde{P}_{\mathbb{C}} / G^{\circ}(\mathbb{Q}) \ \struck{\simeq} \ (P/G^{\circ})_{\mathbb{C}}^{\mathrm{an}} , \]
\struck{et il} et il \ill{} $=$ \ill{} à faire leur restriction à $S$
\struck{\ill{} et} au \ill{}. Les \ill{} que ce sont celles données plus haut …

\struck{N.B. \ill{}}\note{Le N.B.\ qui ferme le feuillet est biffé d'un trait et
n'est plus lisible.}

\page{43}

\begin{enumerate}
  \item[A)] Groupe réductif \struck{$E$} évidemment $G$ sur $\mathbb{Q}$.
  \item[B)] Fibré principal à connexion intégrable $P$ sur $G_{S}$.
  \item[C)] Restriction du groupe $G_{\mathbb{C}}^{\mathrm{an}}$ à $G(\mathbb{Q})$
  dans $P_{\mathbb{C}}^{\mathrm{an}}$, d'où rev.\ \struck{\ill{}}
  $\widetilde{P} \overset{p}{\hookrightarrow} P_{\mathbb{C}}^{\mathrm{an}} \to P$.
  \item[D)] Application holomorphe
  \[ f : \widetilde{P} \longrightarrow \Gamma_{G} \]
  satisfaisant $f(xg) = g^{-1} f(x)$ \ \ ($x \in \widetilde{P}$,
  $g \in G(\mathbb{Q})$) et compatible avec les données B) et $p$ de C).
  \item[E)] Restriction de $\widetilde{P} \times^{G(\mathbb{Q})} G(\mathbb{A})$ à
  $S$, avec compatibilité \struck{\ill{}} verticalement : $G/G^{\circ}$.
\end{enumerate}

\marginal{où $q$, \uncertain{$f'$}}\note{Cette annotation marginale est reliée par
une accolade aux items C), D) et E).}

N.B. La connaissance de A) B) C) et de $f(\xi_{0})$, i.e.\ $\xi_{0} \in \widetilde{P}$
— au-dessus des points génériques de $S$ [\ill{} peut-être \emph{simplement}] —
doit en principe permettre de reconstituer D) tout entier et les données
provenant de E) \ [lorsqu'on sort \ill{} \ill{} que les données provenant d'une
cat.\ de motifs $\mathcal{M}$].

Formule analytique de prolongement de $f$ ?

Relations entre D) et E) ?? \ [\uncertain{Jugendtraum} ?]\note{Le mot est écrit
entre crochets et suivi d'un point d'interrogation ; la lecture est offerte
sous réserve, mais le contexte adélique de E) la rend plausible.}

\marginal{N.B. Quand $f$ est \ill{} sur \uncertain{le pt.\ génér.} de $S$, il
\ill{} de connaître la connexion de $P$, et il suffit de connaître $P_{\xi}$ ?}

\section{Foncteurs pleinement fidèles sur catégories des motifs, liés à théorie de Hodge}

\note{Le titre est celui qu'il inscrit, souligné et précédé d'un $\otimes$
entouré, en tête du feuillet 45. Les six rubriques qui suivent sont numérotées
de sa main 1) à 6) ; chacune est annoncée par une rubrique entre parenthèses —
« Via Hodge », « Via Tate », « Double niveau ».}

\page{45}

Soit $S$ un $\mathbb{Q}$-schéma régulier \add{connexe} muni d'un point
géom.\ $\xi : \operatorname{Spec}(\mathbb{C}) \longrightarrow S$. Le foncteur
$\xi^{*} : \mathrm{Mot}(S) \to \mathrm{Mot}(\mathbb{C})$, $\mathrm{Mot}(S)$
désignant les motifs lisses sur $S$, est fidèle. On regarde
\struck{le foncteur composé} le foncteur composé

\begin{tikzcd}[column sep=small, row sep=small, nodes={font=\scriptsize}]
\mathrm{Mot}(S) \arrow[r, "\xi^{*}"] \arrow[rrr, bend right=25, "(T_{B})_{\xi}"'] & \mathrm{Mot}(\mathbb{C}) \arrow[r, "T_{B+H}"] & \mathrm{Hodge}(\mathbb{Q}, \mathbb{C}) \arrow[r] & \mathrm{Mod}_{f}(\mathbb{Q})
\end{tikzcd}

\note{L'étiquette de la deuxième flèche est réécrite par-dessus une autre ; le
$H$ ajouté après le $B$ est sûr, ce qui est dessous ne l'est pas.}

\marginal{Réseau et vectoriel filtré}

\emph{(Via Hodge)}

1) Supposons $\xi$ localisé au pt.\ générique de $S$. \add{a)} À tout motif $M$
sur $S$, on associe le triple
\[ \{\, T_{DR}(M_{y}),\ T_{B}(M_{\xi}),\ \varphi_{M} \,\} \]
où $\varphi_{M}$ est l'isom.\ can.
\[ \varphi_{M} : T_{B}(M_{\xi}) \otimes_{\mathbb{Q}} \mathbb{C} \ \simeq \ T_{DR}(M_{y}) \otimes_{k(y)} \mathbb{C} . \]
On considère $T_{DR}(M_{y})$ comme \emph{filtré}, la filtration étant telle que
$\varphi_{M}$ la transforme en une filtration \uncertain{déterminée} des
\ill{} de $T_{B}(M_{\xi})$ \ [i.e.\ \ill{}, avec \ill{} conjugué, une bigraduation].

Le foncteur $M \mapsto \ill{}$ ainsi obtenu est pleinement fidèle (moyennant
conj.\ de Hodge simple).

\marginal{Réseau et vectoriel bigradué}

b) Même résultat si on \add{considère} \ill{}
\[ \{\, T_{Hdg}(M),\ T_{B}(M_{\xi}),\ \psi_{M} \,\} \]
$T_{Hdg}(M)$ étant considéré comme \emph{bigradué}, et $\psi_{M}$ étant
l'isomorphisme canonique
\[ \psi_{M} : T_{B}(M_{\xi}) \otimes_{\mathbb{Q}} \mathbb{C} \ \simeq \ T_{Hdg}(M) \otimes_{k(y)} \mathbb{C} . \]

\marginal{Le \struck{foncteur} n'est-il \ill{} pl.\ fidèle si \ill{} la donnée de
la filtration \ill{} $T_{DR}$ ?? \ill{} \,!! de $T_{Hdg}$}

\page{46}

\emph{(Via Tate)}

2) Soit $\pi = \pi_{1}(S, \xi)$, \struck{dont l'un les précédents}. À tout $M$,
associons
\[ \{\, T_{B}(M_{\xi}),\ \rho_{\ell, M} \,\} \]

\note{Un « $+H$ » a été écrit après le $B$ de $T_{B}$, puis biffé : il
écarte ici la structure de Hodge, qu'il rétablira plus bas sous le nom
$T_{B-H}$.}
\struck{i.e.\ $T_{D.H}(M_{\xi})$ est la structure de \struck{Hodge} Hodge
associée, i.e.\ $M_{\xi} \mapsto T_{D,H}(M) = \{ T_{D}(M_{\xi}) + \text{bigraduation} \mapsto T_{D}(M_{y}) \otimes_{\mathbb{Q}} \mathbb{C} \}$,}
\struck{et} où
\[ \rho_{\ell, M} : \pi \longrightarrow \operatorname{Aut}_{\mathbb{Q}_{\ell}}
\bigl( T_{B}(M_{\xi}) \otimes_{\mathbb{Q}} \mathbb{Q}_{\ell} \bigr) . \]

\marginal{$\ell$-Modules galoisiens — réseaux}

Le foncteur ainsi obtenu est pl.\ fidèle si $S$ \struck{un corps $k$
($\subset \mathbb{C}$) \ill{}} est de type fini sur $\mathbb{Q}$, en admettant
comme \struck{Hodge} Tate les conj.\ de \struck{Hodge}, et utilisant rés.\ de
Griffiths.

Si $S$ est loc.\ de t.f.\ sur $\mathbb{C}$, on peut aussi prendre
$\pi' = \pi_{1}(S(\mathbb{C}), \xi)$, et $\rho_{M}$ \add{au lieu de $\rho_{\ell}$},
\[ \rho_{M} : \pi' = \pi_{1}(S(\mathbb{C}), \xi) \longrightarrow
\operatorname{Aut}_{\mathbb{Q}} \bigl( T_{D}(M_{\xi}) \bigr) , \]
à condition de \struck{remplacer} $T_{B}$ par $T_{B-H}$ (structure de Hodge !).

\marginal{ici c'est conj.\ Hodge qu'on utilise !}

Si $S$ est p.\ ex.\ affine, \struck{\ill{}} mais pour une cat.\ de type fini sur
un corps, on a un résultat \emph{analogue}, avec $\pi$ remplacé par un système
projectif de gpes fondamentaux de schémas de t.f.\ sur $\mathbb{Q}$ …

3) À tout $M$, associons
\[ \bigl( T_{DR}(M),\ \sigma,\ T_{B}(M_{y}),\ \varphi \bigr) \]
où $T_{DR}(M)$ est la coh.\ de D.R.\ de $M$ sur $S$, considérée comme
\uncertain{vectoriel} \emph{filtré} \struck{(il suffit de donner la filtration de
$T_{DR}(M)(\xi)$ !)}, $\sigma$ sa stratification (Gauss--Manin) relativement à
\ill{} \add{un corps de base fini alg.\ clos $k = \overline{\mathbb{Q}}$},
$T_{B}(M_{\xi})$ la coh.\ de Betti, et $\varphi$ l'isomorphisme canonique
\[ T_{DR}(M)(\xi) \ \simeq \ T_{B}(M_{y}) \otimes_{\mathbb{Q}} \mathbb{C} . \]
Si $S$ est de \ill{}, on trouve ainsi un foncteur pl.\ fid.\ (moyennant conj.\ de
Hodge, et utilisant Griffiths).

\marginal{\struck{\ill{}} si $k = \mathbb{Q}$, \ill{} n'est-il \ill{} de vecteur la
filtration sur $T_{DR}(M)$ ? \ill{} plus bas.}

\page{47}

\emph{« Double niveau »}

4) Supposons $S = \operatorname{Spec} k$, $k$ \add{une} extension
\emph{algébrique} de $\mathbb{Q}$. Considérons le foncteur qui à $M$ associe
\[ (4.1) \qquad \{\, T_{DR}(M),\ T_{B}(M),\ \varphi \,\} \]
où $\varphi : T_{B}(M) \otimes_{\mathbb{Q}} \mathbb{C} \simeq
T_{DR}(M) \otimes_{k} \mathbb{C}$ ; ici, on ne considère pas $T_{DR}(M)$ comme
filtré, \struck{et} ni $T_{B}(M)$ muni d'une structure de Hodge. Ce foncteur
est-il pl.\ fidèle ??

\struck{Il revient au même de \ill{}} Il est aussitôt ramené au cas où
$k = \overline{\mathbb{Q}}$. La question est alors équivalente à celle si le
foncteur associant à $M$ le système
\[ (4.2) \qquad \{\, T_{B}(M),\ \sigma \,\} \]
où $\sigma$ est le \ill{} canonique sur $T_{B}(M) \otimes_{\mathbb{Q}} \mathbb{C}$,
est \emph{pleinement fidèle} : en effet, $\sigma$ et $T_{DR}(M) \subset
T_{B}(M) \otimes \mathbb{C}$ se déterminent mutuellement
[$T_{DR}(M)$ est l'ens.\ des éléments \ill{} par $\sigma$ …].

Une autre façon de définir \add{$T_{DR} \subset T_{B} \otimes \mathbb{C}$,} cette
donnée, quand on introduit le groupe de Galois
\[ G = \underline{\operatorname{Aut}}^{\otimes}\bigl( T_{B}\struck{(M)} \ (M \mapsto \otimes\,\ill{}) \ \text{de la catégorie } \mathrm{Mot}(k) \bigr) \longrightarrow \]
catégorie de type fini \ill{}, c'est la donnée de
$\dot{g} \in G(\mathbb{C})/G(\overline{\mathbb{Q}})$ ; si $V$ est un
$G$-module (correspondant à un motif $M$), alors $T_{DR}(M)$ n'est autre que le
$\overline{\mathbb{Q}}$-réseau $g(V_{\overline{\mathbb{Q}}}) \subset
V \otimes_{\mathbb{Q}} \mathbb{C}$.

Donc la question est simplement si pour \ill{} $V$, l'inclusion
\[ V^{G} \ \subset \ V \cap g(\overline{V}) \]
est une identité. [On peut se ramener \ill{} que $V$ est un module simple.
[S'il est de dim.\ $1$, c'est une puissance de Tate, et ça marche.]]

\marginal{N.B. $G$ connexe}

\page{48}

\struck{regardons $V$ irréductible (i.e.\ $\overline{V} = V \otimes_{\mathbb{Q}} \overline{\mathbb{Q}}$)}

\[ \struck{\overline{V}^{G} \ \subset \ \overline{V} \cap g(\overline{V})} \]

\struck{et une \ill{} qui est une identité [car $\overline{V}^{G} \cap V = V^{G}$]}.

Or la relation $x \in \overline{V} \cap g(\overline{V})$ équivaut à
\[ x \in \overline{V}, \qquad g^{-1} x \in \overline{V} \]
et implique évidemment, si $\overline{T} \subset \overline{G}$
\struck{le lieu de $g$ sur $G_{\overline{\mathbb{Q}}}$} \struck{les dé\ill{}},
\[ x \in \overline{V} , \quad t\, s^{-1} \in G_{x}(\mathbb{C}) \ \text{ pour } t, s \in \overline{T}(\mathbb{C}) , \]
et \struck{$= \gamma$} équivalant. En d'autres termes, si $G_{x}$ est le
stabilisateur de $x$, qui est un \ill{} agissant dans $G$, la relation
$x \in \overline{V} \cap g(\overline{V})$ équivaut à
$\overline{R} \in G_{x}$ \struck{\ill{}}, i.e.\ $\overline{R}$ est l'adhérence de
$\overline{T} \times \overline{T} \to \overline{G}$, $(s,t) \mapsto s^{-1} t$.
Soit $\overline{R}$ l'image de $\overline{R}$ dans $G$, $H$ le sous-groupe
\struck{\ill{}} $H \subset G$, i.e.\ $H$ \ill{} $x$.

\struck{\ill{}}\note{Suit, encadré, un long passage écrit par-dessus lui-même
puis biffé en travers ; il porte sur le stabilisateur $G_{x}$, sur les invariants
de $G$, sur $\overline{R}$ et sur une condition « il suffit que … pour un
voisinage de \ill{} », et n'est pas récupérable.}

\struck{i.e.\ \ill{} que $\overline{R}$ invariant de \ill{}, $V^{G}$,
i.e.\ $G_{x} = G$ \ill{}, $\overline{R} \subset \ill{}$, $H$ de $G$
(pour \ill{} repr.\ de $G$) \ill{} $H = G$, mais cette condition n'est \ill{}
un $k$-gpe de Borel de $G$ a \ill{} (en \ill{})}

\marginal{La propriété qui \ill{} vraiment \ill{} elle-ci :}

\marginal{N.B. \struck{$G$} $\overline{R}$, $\overline{H}$ ne dépendent pas au
choix de \ill{} quand $G(\overline{\mathbb{Q}})$, \struck{\ill{}}}

\marginal{Regardez notamment le \uncertain{cas des courbes elliptiques}}

5) Bien entendu, dans \ill{} que $k$ est alg.\ sur $\mathbb{Q}$, on peut regarder
le foncteur (4.1). Mais il n'est plus en général pl.\ fidèle,
\add{$\otimes$-cat.} car si $k$ est de t.f.\ sur $\mathbb{Q}$, on
\struck{peut prendre} $M$ \emph{fixé} \add{de} type fini \ill{} de
$\mathrm{Mot}(k)$, si la condition d'engendrer les $k$, \add{on trouve que}
$T_{DR}(M)$ contient $\underline{T_{B}(M)}$, \add{plus haut} \ill{} \ill{} \ill{}
simple \ill{} $T_{B}(M)$ : idiot !

Pour $M$ fixé, on peut espérer tout au moins avoir quelque résultat de
disproportion, de $T_{B}$ et $T_{DR}$ si $k$ est $\pm$ modéré (p.\ ex.\ avec les
répartitions des degrés de transcendance minimum).

\page{49}

Par contre, sans hypothèse sur $k$ sauf d'être \emph{alg.\ clos}, il semble
encore raisonnable d'espérer que (4.2) soit très pleinement fidèle. On est ramené
pour cela au cas $k = \mathbb{C}$. C'est une conséquence de $2$ conjectures
« du double niveau » de 4) $+$ Tate ou Hodge, plus : 6)

\marginal{\uncertain{réseau $+$ connexion}}

6) Si $S$ de type fini \add{\struck{\ill{}} un corps de nbs}, \ill{} : le
\struck{foncteur} de 3), mais sans filtration sur $T_{DR}(M)$. Il résulterait de
la conj.\ du double niveau \struck{qu'elle \add{tout} génér\ill{}} que ce foncteur
est pl.\ fidèle.

[N.B. Si $S$ est le spectre d'un corps de nbs, \uncertain{ceci} se réduit :
l'énoncé du double niveau …)

\section{Remarques sur la caractérisation des classes de coh. D.R. essentiellement algébriques}

\note{Le titre est celui qu'il inscrit, souligné, en tête du feuillet 50 ; la
seconde moitié en est « par bigraduation ou par \uncertain{vérité} … », dont le
dernier mot n'est pas assuré. Le mot « essentiellement » est une addition
interlinéaire au-dessus d'un mot biffé.}

\page{50}

$K \subset \mathbb{C}$ sous-corps, $M$ \add{motif sur $K$}, d'où $T_{DR}(M)$.
Considérons $T_{DR}(M)$ \add{$= \operatorname{Hom}(1, M) \otimes_{\mathbb{Q}} K$} ;
on voit que \ill{} un \ill{}-module de $T_{DR}(M)$, \struck{et qui} \ill{} :

\begin{enumerate}
  \item[a)] contenu dans $T_{DR}(M)^{00}_{\mathbb{C}}$ ;
  \item[b)] \struck{contenu dans $\operatorname{Hom}(1, M) \otimes_{\mathbb{Q}} (K \cap \mathbb{R}) \to$}
  contenu dans $\bigl( T_{DR}(M)_{\mathbb{C}} \bigr)_{\mathbb{R}}$
  \struck{\add{et $=$} $\operatorname{Hom}(1, M) \otimes_{\mathbb{Q}} (K \cap \mathbb{R}) \to$ contenu dans $\bigl( T_{DR}(M)_{\mathbb{C}} \bigr)_{\mathbb{R}}$}.
\end{enumerate}

A-t-on une réciproque ?

1°) Pas pour \struck{b)} \add{a) et b) simultanément} \struck{\ill{}}
\add{pour $2$ cat.\ de motifs $M$ \ill{}} \ill{} $G$ \ill{}.

[Car des motifs elliptiques \ill{} \ill{} donc car
$T_{DR}(M)^{00}_{\mathbb{C}} = T_{DR}(M)_{\mathbb{C}}$, (d'autre part
\ill{} : b) \ill{} mais ses l'\ill{} \ill{} galoisien \ill{} de gpe $G$ de
$\mathbb{Q}$ \ill{} pour $M \to$ \emph{tout vide} \ill{}
$M = \mathbb{Q}[\mathbb{C}]$, qui donne $T_{DR}(M) = \bot$ …] ; donc dans ce cas
on \ill{} réciproque : a) et b) ).

En effet, si $G$ est \struck{\ill{}} \ill{} de poids $0$ ], la condition a) est
vide, donc réc.\ de a) vraie ssi $G = (e)$ ; \struck{b)} la réciproque
\ill{} (d'autre part \ill{} galoisien \ill{} de gpe $G$ de $\mathbb{Q}$ \ill{}
\emph{tout vide} [i.e.\ \ill{} avec le $G$-module \ill{} dans ce cas on
\ill{} a) et b) )].

Il faut donc \ill{} \ill{} $G$ connexe, on \ill{} p.\ ex.\ $K = \overline{\mathbb{Q}}$
\ill{} \ill{}. On peut de plus supposer la condition, \ill{} \add{immersions}
$K \hookrightarrow \mathbb{C}$.

2°) La réciproque \ill{} a) $\to$ \ill{} \ill{} \ill{} ($=$ pour
$\overline{\mathbb{Q}}$), il s'agit de prendre \ill{} \ill{} gpe de Galois
motivique $G$ \ill{} de dim.\ $\geqslant 3$. On note qu'en \ill{} fixant la
$(K \hookrightarrow \mathbb{C})$, \ill{} \ill{} que $G$
\add{$= \operatorname{ad}(P)$, un $G$ commut.} (qui définirent) (la bigraduation,
\[ j_{1}, j_{2} : G_{K} \overset{\rho}{\rightrightarrows} G \]
\[ V_{\mathbb{C}} \ \simeq \ V' \otimes_{K} \mathbb{C} \qquad
[\, V' = P \times^{G_{K}} V, \quad
\rho = \struck{\underline{\operatorname{Isom}}^{\otimes}} \bigl( T_{B} \otimes k,\ T_{DR} \bigr) \,] \]
Donc une représentation \emph{fidèle} de gpe \ill{} $0$ de $G$ donnera,
\struck{\ill{}} $T$ l'image de \add{$(j_{1}, j_{2}) : G_{K} \times G_{K} \to G'$},
\ill{} $V'$ \ill{} pour laquelle
\[ {V'}^{\,T} \ \bigl\{ = T_{DR} \cap (T_{DR})^{00}_{\mathbb{C}} \bigr\}
\ \ill{} \ \neq \ {V'}^{\,G'} . \]

\note{Le feuillet 50 est écrit par-dessus lui-même : la marge de gauche porte, au
niveau de 2°), un bloc de trois ou quatre lignes entièrement recouvert de
ratures, dont rien ne se lit.}

\page{51}

3°) Reste la question d'une réciproque à b) si $K$ \ill{}, en prenant
\add{au moins} (si $K \neq \overline{\mathbb{Q}}$) \emph{toutes} les
imm.\ $K \to \mathbb{C}$, et en \emph{conjuguant} aussi au besoin avec la condition a).

\note{Le feuillet s'arrête sur ces trois lignes, et la question 3°) reste posée :
le dossier n'a pas de suite.}

\end{document}
